\chapter{Performance Metrics}

\begin{ppnote}
These notes are from Lecture~21 on performance metrics for parallel programs.
They follow the order in which the material was presented in class, fix a
few mistakes in the printed slides (flagged in pitfall boxes), and fill in the
derivations written out on the board but not on the deck. The closing
section on iso-efficiency and the ``strongly~/~weakly scalable'' classification
was given verbally and does not appear on the slides at all, so it has been
written up from notes.\\[2pt]
\textit{Lecturer:} Prof. Subodh Sharma, IIT Delhi (COL380, Introduction to
Parallel \& Distributed Programming).\\
\textit{Transcribed by:} Anirudha Saraf, Manvendra Rajpurohit, Gaurav Raju Shende.
\end{ppnote}

% ==========================================================
\section{Why we care about performance metrics}
% ==========================================================

Correctness is one half of the job. The whole point of writing a parallel
program in the first place is to finish a workload faster, or to handle a
bigger workload than a single machine can, or to use the hardware more
economically. None of those claims mean anything unless we have a way to
measure them.

The point was made at the start of class that this material could really
have been covered on day one --- the moment you write your first OpenMP, MPI
or CUDA program, you start wanting to know whether the work you put in
actually paid off. The answer has several parts: how does the program
behave as you give it more cores, how does it behave as the problem itself
grows, and how much of the runtime is being eaten by overheads
(communication, cache coherence, synchronisation, ordering fences, and so on).

This chapter pins down the vocabulary --- speedup, cost, overhead,
efficiency, and scalability --- and then presents the two laws (Amdahl and
Gustafson) that say something about how far parallelism can take you. The
closing topic, iso-efficiency, ties everything together by asking how fast
the problem must grow if we want to keep the processors usefully busy.

\begin{ppkey}
Given a sequential program and $p$ processors: how much faster does it run,
how well are we using those $p$ processors, and how big does the problem
need to be for that quality to hold up?
\end{ppkey}

% ==========================================================
\section{The RAM model}
% ==========================================================

Before defining anything, we fix a model of computation. Throughout this
course we use the standard \ppterm{RAM (Random Access Machine)} model, with
these assumptions:

\begin{itemize}
  \item Every simple arithmetic or logical operation takes one unit of time.
  \item Every local memory access (by a thread or process) takes one unit of time.
  \item Each thread or process is an independent unit of execution with its
        own local view of memory.
  \item If process $i$ sends a constant-size message to process $j$, the
        transmission also takes one unit of time.
\end{itemize}

The model is an idealisation --- real hardware has caches, NUMA effects,
network contention, and so on --- but it is the baseline. The
communication-aware extensions later in the chapter refine it without
throwing it away.

% ==========================================================
\section{The basic definitions}
% ==========================================================

\subsection{Sequential and parallel time}

Throughout we write
\begin{itemize}
  \item $T_1$ for the wall-clock time of the best sequential algorithm for
        the problem on a single processor, and
  \item $T_p$ for the wall-clock time of the parallel algorithm on $p$ processors.
\end{itemize}
A small but important point: $T_1$ is the time of the
\emph{fastest known serial algorithm}, not the parallel algorithm running on
a single processor. Otherwise we could make a slow parallel algorithm look
good just by comparing it to itself.

\subsection{Speedup}

\begin{ppdefn}
\textbf{Speedup.}
\[
S_p \;:=\; \frac{T_1}{T_p}.
\]
$S_p$ is how many times faster the parallel run was, compared to the serial baseline.
\end{ppdefn}

There are three regimes worth knowing:

\begin{itemize}
  \item \ppemph{Linear (ideal) speedup:} $S_p = p$. Throwing $p$ resources
        at the problem reduced the time by exactly a factor of $p$.
  \item \ppemph{Sub-linear speedup:} $S_p < p$. The usual case. The typical
        culprits are communication latency, synchronisation, cache coherence
        traffic, thread management, memory management, and the ``hundred
        small things'' that seep into the mixture.
  \item \ppemph{Super-linear speedup:} $S_p > p$. Rare but real. Usually it
        is a cache effect (the per-processor working set now fits in L2) or a
        search-parallelism effect (parallel exploration finds the answer earlier).
\end{itemize}

\subsection{Cost}

\begin{ppdefn}
\textbf{Cost.}
\[
C_p \;=\; p \times T_p.
\]
The total computational effort: processor-seconds spent.
\end{ppdefn}

Speedup alone is dangerous. The slide gives a nice example: $T_1 = 100$\,s
and $T_8 = 20$\,s gives $S_8 = 5$, which sounds great. But the cost is
$C_8 = 8 \times 20 = 160$ processor-seconds --- 60 more than the serial run
did with one processor. We got the answer faster, but we did $60\%$ more
total work to get it. If energy were our metric (\ppterm{flops per joule},
the modern figure of merit for \emph{green supercomputers}), this parallel
run is actually worse.

\subsubsection*{Cost-optimality}
A parallel algorithm is \ppterm{cost-optimal} when
\[
p \cdot T_p \;=\; \mathcal{O}(T_1),
\]
i.e.\ the parallel total work is asymptotically the same as the sequential
work. A rather brutal counter-example was given in class to make the point:

\begin{ppexample}
\textbf{Why cost matters: matrix multiplication.}
Take $n \times n$ matrix multiplication, where the serial algorithm takes
$T_1 = \Theta(n^{3})$. Suppose we throw $p = n^{2}$ processors at it and
somehow get $T_p = n$. Then
\[
S_p = \frac{n^{3}}{n} = n^{2} \quad\text{(huge speedup),}
\qquad
C_p = n^{2}\cdot n = n^{3}.
\]
This is fine --- cost-optimal. But suppose instead $T_p = n^{2}$:
\[
S_p = n, \qquad C_p = n^{2}\cdot n^{2} = n^{4}.
\]
We pay $n^{4}$ total work for a speedup of only $n$. The cost has gone from
$n^{3}$ to $n^{4}$, which is one full factor of $n$ of wasted work. The
parallel solution is \ppterm{unscalable} --- the speedup looks fine in
isolation but the bookkeeping reveals the trap.
\end{ppexample}

\subsection{Parallelism overhead}

\begin{ppdefn}
\textbf{Total overhead.}
\[
T_o \;=\; C_p - T_1 \;=\; p \cdot T_p - T_1.
\]
\end{ppdefn}

This is the part of the parallel work that is not the serial algorithm's
essential computation. It bundles together:
\begin{itemize}
  \item communication between processors (message passing, cache coherence
        traffic),
  \item synchronisation (barriers, locks, atomic updates, ordering fences),
  \item idle time from load imbalance,
  \item thread~/~process management (creation, scheduling, teardown),
  \item redundant computation done to avoid communication.
\end{itemize}

Cost-optimality is essentially the statement that the overhead does not
grow faster than $T_1$: $T_o/T_1 \to 0$ as the problem and $p$ scale together.
If $T_o$ grows faster than $p$, scaling fails entirely --- adding processors
only makes things worse.

\subsection{Efficiency}

\begin{ppdefn}
\textbf{Efficiency.}
\[
E_p \;=\; \frac{T_1}{C_p} \;=\; \frac{T_1}{p \cdot T_p} \;=\; \frac{S_p}{p}.
\]
The fraction of total processor-time spent doing useful work. $E_p$ lies in
$(0,1]$, and $E_p = 1$ is the ideal case.
\end{ppdefn}

The slide example is worth pausing on:

\begin{ppexample}
\textbf{Speedup hides what efficiency reveals.}
\begin{itemize}
  \item $S_{8} = 7$ gives $E_{8} = 7/8 = 0.875$. Processors are
        $87.5\%$ utilised --- excellent.
  \item $S_{64} = 20$ gives $E_{64} = 20/64 \approx 0.3125$. On 64 processors,
        only $31\%$ of the aggregate capacity is doing useful work.
\end{itemize}
$S_{64}$ is almost three times $S_8$, but the 8-processor configuration is
twice as efficient. Speedup numbers reported on their own are half-truths;
always look at efficiency next to them.

The natural follow-up question --- \emph{if we want to keep efficiency at
$0.875$ even on 64 processors, how much do we need to grow the problem?} ---
is what iso-efficiency answers in
Section~\ref{sec:isoefficiency}.
\end{ppexample}

% ==========================================================
\section{Scalability}
\label{sec:scalability}
% ==========================================================

A speedup number is reported for one specific (problem size, processor
count) pair, which is why it is so often misleading. Scalability asks how
the metrics \emph{evolve} when we change either knob, or both.

From here on we treat $T$, $S$, $E$, $C$ as functions of two variables: the
problem size $n$ and the processor count $p$.

The reason this matters is simple. If we hold $n$ fixed and keep increasing
$p$, the per-processor work shrinks, but the communication and
synchronisation costs do not, so the speedup curve flattens out. If we
instead grow $n$ together with $p$, the parallel \emph{opportunity}
grows --- there is more work to spread across the processors --- while the
serial portion of the program stays roughly the same size in absolute terms.
The two situations behave very differently, which is why they get separate
names.

\subsection{Strong scaling}

\begin{ppdefn}
\textbf{Strong scaling.}
Hold $n$ fixed and let $p$ grow. Study how $T(n,p)$ behaves.
Ideal strong scaling:
\[
T(n,p) \;=\; \frac{T(n,1)}{p}, \qquad \text{equivalently } S_p = p.
\]
\end{ppdefn}

Strong scaling rarely holds in practice for large $p$. Once the
per-processor work shrinks enough, the overheads dominate. A useful way to
write the parallel runtime is
\[
T(n,p) \;=\; T_{\text{comp}}(n,p) \;+\; T_{\text{comm}}(n,p) \;+\; T_{\text{sync}}(n,p) \;+\; T_{\text{idle}}(n,p).
\]
For fixed $n$, $T_{\text{comp}}$ falls like $1/p$, but the others typically
stay constant or grow with $p$, so eventually the second group dominates the
first. On top of this, load imbalance gets harder to avoid as the chunks get
smaller, so beyond some point adding more processors actively slows the
program down.

\begin{pppitfall}
\textbf{Slide correction.} The slide writes ``Strong scaling breaks down:
$T(n,p) = T_{\text{comp}} + T_{\text{comm}} + T_{\text{sync}} + T_{\text{idle}}$.''
This was corrected in class. The decomposition is the \emph{general} formula
for parallel runtime --- it is always true. What actually causes strong
scaling to fail is that the \emph{overhead terms}
($T_{\text{comm}} + T_{\text{sync}} + T_{\text{idle}}$) grow with $p$.
Without a lower bound on those, throwing more processors stops helping.
\end{pppitfall}

\subsection{Weak scaling}

\begin{ppdefn}
\textbf{Weak scaling.}
Each processor is given a fixed workload $m$, so the total problem size is
$n = mp$. Ideal weak scaling:
\[
T(mp,\,p) \;=\; T(m,\,1),
\]
i.e.\ runtime stays constant as $p$ and $n$ grow together.
\end{ppdefn}

\begin{pppitfall}
\textbf{Slide correction.} The slide states ideal weak scaling as
$T(mp, p) = T(m,1)/p$. This is wrong --- that formula is what we would
write for ideal \emph{strong} scaling on a problem of size $m$. The whole
point of weak scaling is that each processor's local work stays the same,
so the wall-clock time should stay constant, not shrink. The correct
formula is
\[
T(mp,p) = T(m,1) \quad\text{(or at worst approximately constant in $p$).}
\]
\end{pppitfall}

\begin{ppexample}
\textbf{Slide example.}
\begin{center}
\begin{tabular}{@{}ccc@{}}
\toprule
$p$ & Total problem size & Time (s) \\ \midrule
1  & 1\,000\,000 & 10 \\
2  & 2\,000\,000 & 11 \\
4  & 4\,000\,000 & 13 \\ \bottomrule
\end{tabular}
\end{center}
Ideal weak scaling would have held the time at $10$\,s. The drift from $10$
to $11$ to $13$\,s is real communication overhead, but it is small and the
growth is far sub-linear in $p$. This is the kind of ``approximately
constant'' behaviour described as acceptable in class.
\end{ppexample}

% ==========================================================
\section{Amdahl's law}
% ==========================================================

The setup: split the sequential time $T_1$ into
\begin{itemize}
  \item a fraction $f$ that is inherently sequential and cannot be parallelised, and
  \item a fraction $(1-f)$ that can be parallelised perfectly.
\end{itemize}

On $p$ processors, the sequential part still costs $f \cdot T_1$, while the
parallel part costs $(1-f)\cdot T_1 / p$. Adding them up:

\begin{ppdefn}
\textbf{Amdahl's law.}
\[
T_p \;=\; T_1\!\left(f + \frac{1-f}{p}\right),
\qquad
S_p \;=\; \dfrac{1}{f + \tfrac{1-f}{p}},
\]
and as $p \to \infty$:
\[
S_{\infty} \;=\; \frac{1}{f}.
\]
\end{ppdefn}

\begin{pppitfall}
\textbf{Slide correction.} The slide writes
\[
T_p \;=\; T_1\bigl(1 + 1 - f/p\bigr),
\]
which is just garbled --- the parentheses and signs do not mean what they
should. The correct expression is
\[
T_p \;=\; T_1\!\left(f + \frac{1-f}{p}\right),
\]
which is consistent with the (correct) $S_p = 1/(f + (1-f)/p)$ printed
right below it.
\end{pppitfall}

\subsection{What Amdahl's law actually tells us}

For a fixed-size problem, the serial fraction sets a hard ceiling on how
fast you can ever go, no matter how much hardware you throw at it.

\begin{itemize}
  \item $f = 0.10$ ($10\%$ serial) gives $S_\infty = 10$. This is the example on the slide.
  \item $f = 0.01$ ($1\%$ serial) gives $S_\infty = 100$.
  \item $f = 0.001$ ($0.1\%$ serial) gives $S_\infty = 1000$.
\end{itemize}

The lesson is that the serial fraction governs everything. Speeding up
code that is already parallel has diminishing returns; even a tiny
reduction in the serial portion can raise the ceiling dramatically.

\begin{ppkey}
\textbf{The NVIDIA example given in class.}
Imagine you are a product manager at NVIDIA with a \$10\,M budget for next
quarter and a known workload whose serial fraction you can estimate. The
capacity-planning calculation \emph{is} Amdahl's law: compute
$S_\infty = 1/f$, identify the $p$ at which the curve plateaus, and that is
roughly the maximum number of GPUs worth buying. Beyond that, you are
paying for cards that cannot make the job any faster.
\end{ppkey}

\subsection{What the speedup curves look like}

\begin{figure}[h]
\centering
\begin{tikzpicture}
\begin{axis}[
  width=0.85\linewidth, height=7cm,
  xlabel={Number of processors $p$},
  ylabel={Speedup $S_p$},
  xmin=1, xmax=64, ymin=0, ymax=22,
  grid=both, grid style={gray!20},
  legend pos=north west,
  legend style={font=\small},
  every axis plot/.append style={thick},
  title={Amdahl's law for several values of $f$}
]
\addplot [domain=1:64, samples=200, color=PPBlue]
  {1/(0.05 + 0.95/x)};
\addlegendentry{$f=0.05$ \;($S_\infty=20$)}

\addplot [domain=1:64, samples=200, color=PPWarn]
  {1/(0.10 + 0.90/x)};
\addlegendentry{$f=0.10$ \;($S_\infty=10$)}

\addplot [domain=1:64, samples=200, color=PPGood]
  {1/(0.25 + 0.75/x)};
\addlegendentry{$f=0.25$ \;($S_\infty=4$)}

\addplot [domain=1:64, samples=200, color=PPGray]
  {1/(0.50 + 0.50/x)};
\addlegendentry{$f=0.50$ \;($S_\infty=2$)}
\end{axis}
\end{tikzpicture}
\caption{Each curve plateaus at $1/f$ no matter how many processors are added.}
\label{fig:amdahl-curves}
\end{figure}

% ==========================================================
\section{Gustafson's law}
% ==========================================================

Amdahl's law assumes the problem size is fixed, which is rarely how things
work in practice. Most of the time we buy more processors precisely because
we want to solve \emph{bigger} problems --- not the same one slightly
faster. Gustafson noticed that in this scaled-workload setting, the serial
portion of a program tends to stay roughly the same in absolute terms,
while the parallel portion grows with $p$. So the effective serial
fraction \emph{shrinks} as we scale up, and the ceiling becomes much more
forgiving.

\subsection{Setup}

Take the $p$-processor run as the reference, with runtime $T(n,p)$, and
split it as
\begin{itemize}
  \item $f \cdot T(n,p)$ for the sequential portion,
  \item $(1-f)\cdot T(n,p)$ for the parallel portion.
\end{itemize}

If we ran the same workload on a single processor, the parallel portion
would take $p$ times longer (one processor doing what $p$ did) and the
serial portion would be unchanged. So
\[
T(n,1) \;=\; f\cdot T(n,p) \;+\; p \cdot (1-f) \cdot T(n,p).
\]

\begin{ppdefn}
\textbf{Gustafson's law.}
\[
S_p \;=\; \frac{T(n,1)}{T(n,p)} \;=\; f + (1-f)\cdot p \;=\; p - f(p-1).
\]
\end{ppdefn}

This grows linearly in $p$, with no $1/f$ ceiling.

\begin{pppitfall}
\textbf{The $f$ in Gustafson is not the $f$ in Amdahl.}
The slide reuses the symbol $f$ for both laws, which trips students up
every year and was the subject of a long discussion in this lecture. In
Amdahl's law, $f$ is the serial fraction of the \emph{serial} runtime
$T_1$. In Gustafson's law, $f$ is the serial fraction of the
\emph{parallel} runtime $T(n,p)$. They are different numbers with
different physical meanings, and treating them as the same is the source of
the natural objection in the next subsection.

Some textbooks rename Gustafson's $f$ to $s$ to avoid the collision; we
will mention this where it helps.
\end{pppitfall}

\subsection{If Amdahl limits speedup, how does Gustafson avoid it?}
\label{sec:reconcile}

A natural question, asked by Ambar in class: if Amdahl says the speedup is
capped at $1/f$, how can Gustafson promise unbounded speedup for the same
program?

\begin{ppnote}
\textbf{Reconciling the two laws.}
If the two $f$'s were really the same number, equating the two speedup
formulas would force
\[
\frac{1}{f + (1-f)/p} \;\stackrel{?}{=}\; f + (1-f)p,
\]
which is plainly false --- the left side is bounded by $1/f$, the right
side grows without bound.

The way out is that Gustafson's $f$ is the serial fraction of the runtime
on $p$ processors, and that fraction \emph{depends on the problem size}
even though $f$ doesn't have $n$ written in it. As $n$ grows, the
sequential portion stays roughly the same in absolute time while the
parallel portion grows, so the proportion $f$ shrinks. Amdahl's $f$ does
not shrink, because it is tied to the fixed-size serial run.

So the two laws answer two different questions:
\begin{itemize}
  \item \textit{Amdahl}: this exact problem, more processors --- how fast?
  \item \textit{Gustafson}: more processors and a proportionally bigger
        problem --- how does the bigger problem on $p$ processors compare
        to the same bigger problem on one processor?
\end{itemize}
\end{ppnote}

\subsection{Reading Gustafson's law}

\begin{itemize}
  \item If $f = 0.05$ on the parallel run and $p = 100$, Gustafson gives
        $S_p = 0.05 + 0.95 \times 100 = 95.05$. Amdahl with the same
        nominal $f = 0.05$ would have capped this at $1/0.05 = 20$.
  \item The form $S_p = p - f(p-1)$ is illuminating: the speedup is
        ``almost $p$'' minus a small correction. As $p$ grows, the
        correction is dwarfed by the $p$ and the speedup keeps climbing.
\end{itemize}

\subsection{Side-by-side comparison}

\begin{table}[h]
\centering
\renewcommand{\arraystretch}{1.3}
\begin{tabular}{@{}p{0.22\linewidth} p{0.34\linewidth} p{0.34\linewidth}@{}}
\toprule
 & \textbf{Amdahl} & \textbf{Gustafson} \\ \midrule
Problem size & Fixed & Scales with $p$ \\
Meaning of $f$ & Serial fraction of $T_1$ & Serial fraction of $T(n,p)$ \\
Depends on problem size? & No & Yes (implicitly, through $T(n,p)$) \\
Speedup formula & $S_p = \dfrac{1}{f + (1-f)/p}$ & $S_p = f + (1-f)p$ \\
Limit as $p\to\infty$ & $1/f$ & Linear in $p$, no ceiling \\
Governs & Strong scaling & Weak scaling \\
\bottomrule
\end{tabular}
\caption{Amdahl vs.\ Gustafson at a glance.}
\label{tab:amdahl-vs-gustafson}
\end{table}

\begin{figure}[h]
\centering
\begin{tikzpicture}
\begin{axis}[
  width=0.85\linewidth, height=7cm,
  xlabel={Number of processors $p$},
  ylabel={Speedup $S_p$},
  xmin=1, xmax=64, ymin=0, ymax=65,
  grid=both, grid style={gray!20},
  legend pos=north west,
  legend style={font=\small},
  every axis plot/.append style={thick},
  title={Amdahl vs.\ Gustafson at $5\%$ serial fraction}
]
\addplot [domain=1:64, samples=200, color=PPWarn]
  {1/(0.05 + 0.95/x)};
\addlegendentry{Amdahl ($f=0.05$): ceiling $=20$}

\addplot [domain=1:64, samples=200, color=PPGood]
  {0.05 + 0.95*x};
\addlegendentry{Gustafson ($f=0.05$): linear}

\addplot [domain=1:64, samples=2, color=PPGray, dashed]
  {x};
\addlegendentry{Ideal $S_p = p$}
\end{axis}
\end{tikzpicture}
\caption{Same nominal $5\%$ serial fraction, two very different
predictions. They disagree because they are not really talking about the
same quantity (see Section~\ref{sec:reconcile}).}
\label{fig:amdahl-vs-gustafson}
\end{figure}

% ==========================================================
\section{Iso-efficiency: how fast must the problem grow?}
\label{sec:isoefficiency}
% ==========================================================

The lecture closed by going back to the question that came up when we
first defined efficiency. We saw that $E_{64} = 0.3125$ was much worse
than $E_8 = 0.875$ for the same program. The natural inverse question is:

\begin{ppkey}
\textbf{The iso-efficiency question.}
At what rate must the problem size $n$ grow with the processor count $p$
so that $E(n, p)$ stays at some chosen constant $k$?
\end{ppkey}

This is the iso-efficiency view. It pulls together everything we have
done: efficiency is the quality, $n$ and $p$ are the controls, and we ask
how the controls must move together.

\subsection{Derivation}

Start from the definition,
\[
E(n, p) \;=\; \frac{T(n,1)}{p \cdot T(n,p)}.
\]
Use $T_o = p \cdot T(n,p) - T(n,1)$ to rewrite
$T(n,p) = (T(n,1) + T_o)/p$, which gives
\[
E(n,p) \;=\; \frac{1}{1 + T_o(n,p)/T(n,1)}.
\]
If we want this to equal a target efficiency $k < 1$, we solve for $T(n,1)$:
\[
T(n,1) \;=\; \frac{k}{1-k}\cdot T_o(n,p).
\]
$T(n,1)$ is essentially the work $W(n)$ done by the serial algorithm, so
we can write
\[
W(p) \;=\; \frac{k}{1-k}\cdot T_o(n,p).
\]
This is the \ppterm{iso-efficiency function}: it tells us how much work
the problem must contain to absorb the overhead and keep $E$ at $k$.

\subsection{Strongly scalable vs.\ weakly scalable}

The growth rate of $W(p)$ summarises everything about an algorithm's
scalability. The classification given verbally in class:

\begin{ppdefn}
\textbf{Strongly vs.\ weakly scalable.}
Let $W(p)$ be the problem size needed to hold efficiency at a chosen $k$.
\begin{itemize}
  \item If $W(p)$ grows \emph{slowly} with $p$ --- constant, logarithmic,
        or at worst a modest polynomial --- the algorithm is
        \ppterm{strongly scalable}.
  \item If $W(p)$ has to grow as fast as $p$ (or faster) just to hold
        efficiency steady, the algorithm is only \ppterm{weakly scalable}.
\end{itemize}
\end{ppdefn}

\begin{pppitfall}
\textbf{Don't confuse the two ``strong~/~weak'' meanings.}
``Strong scaling'' and ``weak scaling''
(Section~\ref{sec:scalability}) describe \emph{experimental regimes} ---
whether you hold $n$ fixed or grow it with $p$ when you run your benchmark.
``Strongly scalable'' and ``weakly scalable'' here describe an
\emph{algorithmic property} --- how fast its iso-efficiency function grows.
Same words, different things; pay attention to context. Different
textbooks also use ``scalable parallel algorithm'' in slightly different
ways. The iso-efficiency function itself is unambiguous, so when in doubt,
just write it down.
\end{pppitfall}

\subsection{Some rules of thumb}

\begin{itemize}
  \item $W(p) = \Theta(p)$ --- good. Doubling processors needs doubling
        the work to keep efficiency constant. Many dense linear algebra
        routines sit here.
  \item $W(p) = \Theta(p \log p)$ --- still respectable. Typical of
        parallel sorts and prefix-sum algorithms.
  \item $W(p) = \Theta(p^2)$ or worse --- poor. The problem has to blow
        up quadratically just to stand still on efficiency. Algorithms in
        this category are usually redesigned or avoided at large $p$.
\end{itemize}

This also closes the loop on the example from earlier: if we want the
$64$-processor run to keep the $E_8 = 0.875$ efficiency we measured at
$p=8$, the iso-efficiency function tells us exactly how much we need to
grow $n$. A library writer who hands users a fixed implementation without
telling them this growth rate is, as it was put in class, selling a
``facade'' of parallelism.

% ==========================================================
\section{A few practical points}
% ==========================================================

The two classical laws are clean but they paper over a few things that
matter in real systems. None of these were on the slides; they are
standard things worth keeping in mind.

\subsection{Communication-aware models}

Both Amdahl and Gustafson assume the parallel portion scales perfectly.
Real code does not. A simple refinement is
\[
T(n,p) \;=\; \frac{W(n)}{p} \;+\; T_{\text{comm}}(n,p),
\]
where $T_{\text{comm}}$ is typically of the form $\alpha + \beta \cdot m$
for a message of $m$ bytes, $\alpha$ being the latency and $\beta$ the
inverse bandwidth. In many distributed settings $T_{\text{comm}}$ grows
like $\log p$ or $\sqrt{p}$, which slowly eats into the speedup as $p$
grows.

\subsection{Super-linear speedup is real}

The textbook definitions rule out $S_p > p$, but it does happen, usually
for one of three reasons:
\begin{itemize}
  \item \ppemph{Cache:} the partitioned working set fits in faster memory
        than the original did.
  \item \ppemph{Search:} parallel branch-and-bound prunes the tree earlier.
  \item \ppemph{Memory bandwidth:} more nodes mean more aggregate memory
        bandwidth, not just more cores.
\end{itemize}

\subsection{Energy and the green supercomputer}

The rise of the ``green supercomputer'' movement was mentioned in class.
Once GPGPUs entered HPC, raw FLOPS stopped being a sufficient figure of
merit --- two systems with the same FLOPS could burn very different
amounts of electricity. The metric people switched to is
\ppterm{flops per joule}: useful computation per unit of energy. Looked at
this way, the cost example we did earlier ($C_p = 160$ vs.\ $T_1 = 100$)
is not just a wall-clock story but a $60\%$ extra energy bill for the same
answer.

\subsection{What to actually report}

\begin{ppkey}
\textbf{What to report when benchmarking a parallel program.}
\begin{enumerate}
  \item $T_1$ (best serial) and $T_p$ for several $p$.
  \item Both strong-scaling and weak-scaling curves where they make sense.
  \item Efficiency $E_p$, not just speedup --- it is where the wasted
        work shows up.
  \item A breakdown of $T_p$ into computation, communication,
        synchronisation, idle.
  \item The iso-efficiency function (or an empirical estimate) if the
        algorithm is intended for large $p$.
  \item Problem size: without $n$, the numbers don't mean anything.
\end{enumerate}
\end{ppkey}

% ==========================================================
\section{Summary}
% ==========================================================

\begin{table}[h]
\centering
\renewcommand{\arraystretch}{1.3}
\begin{tabular}{@{}l l l@{}}
\toprule
\textbf{Quantity} & \textbf{Definition} & \textbf{Ideal value} \\ \midrule
Speedup $S_p$ & $T_1 / T_p$ & $p$ \\
Cost $C_p$ & $p \cdot T_p$ & $\Theta(T_1)$ \\
Overhead $T_o$ & $C_p - T_1$ & $\to 0$ asymptotically \\
Efficiency $E_p$ & $T_1 / (p \cdot T_p) = S_p / p$ & $1$ \\
Strong scaling & $T(n,p)$ for fixed $n$ & $T(n,1)/p$ \\
Weak scaling & $T(mp,p)$ for fixed $m$ & $T(m,1)$ \\
Amdahl & $S_p = 1/(f + (1-f)/p)$ & $\to 1/f$ \\
Gustafson & $S_p = f + (1-f)p$ & linear in $p$ \\
Iso-efficiency $W(p)$ & work needed to hold $E = k$ & slow-growing in $p$ \\
\bottomrule
\end{tabular}
\caption{Summary of the metrics introduced in this chapter.}
\label{tab:metrics-summary}
\end{table}

\begin{ppkey}
\textbf{Things to take away.}
\begin{enumerate}
  \item Speedup without efficiency lies. $S_{64} = 20$ sounds great until
        you notice $E_{64} = 0.31$.
  \item Strong and weak scaling are different experiments and tell
        different stories. Report both when both make sense.
  \item Amdahl and Gustafson don't actually contradict each other --- they
        are answering different questions, and their $f$'s mean different
        things.
  \item The serial fraction is the enemy. Tiny serial sections dominate
        at large $p$.
  \item Overheads are not a footnote. Communication, synchronisation, and
        idle time are where most speedups die.
  \item Iso-efficiency is the honest scalability metric --- it forces you
        to say ``how fast at what size?''.
\end{enumerate}
\end{ppkey}

\section{Exercises}
\begin{ppexercises}
  \ppexercise{A program has serial fraction $f = 0.02$. What is the
  maximum speedup achievable under Amdahl's law? At what $p$ does the
  speedup curve reach $90\%$ of this ceiling?}

  \ppexercise{For matrix multiplication with $T_1 = \Theta(n^3)$ and a
  parallel implementation achieving $T_p = n^2 / \sqrt{p} + \log p$,
  derive the iso-efficiency function and classify the algorithm as
  strongly or weakly scalable.}

  \ppexercise{Repeat the cost analysis of the matrix-multiplication
  example with $T_p = n \log n$ on $p = n^2$ processors. Is the algorithm
  cost-optimal? Compute $S_p$, $E_p$, and $T_o$.}

  \ppexercise{Show algebraically that Amdahl's $S_p$ and Gustafson's
  $S_p$ are equal at $p = 1$, and explain what that means physically.}

  \ppexercise{A weak-scaling experiment yields times $T(mp, p)$ that grow
  like $\log p$. Is this acceptable? Justify in terms of total work and
  efficiency.}
\end{ppexercises}
